% !TEX root = ../main.tex
\section{Setting: value functions for reward-guided diffusion}
\label{sec:02_setting}

% ---------------------------------------------------------------------------
\subsection{The base process and the tilted target}

Fix a time horizon $t\in[0,1]$ and a base diffusion on $\mathbb{R}^{d}$,
\begin{equation}
  \mathrm{d}X_t = f(X_t,t)\,\mathrm{d}t + \sqrt{2a}\,\mathrm{d}W_t,
  \qquad X_0 = 0,
  \label{eq:base-sde}
\end{equation}
with $a>0$ a scalar diffusion coefficient and $W$ a standard $d$-dimensional
Brownian motion. Write $\mathbb{P}$ for the induced path measure, $\E_{\Base}$
and $p_{\Base}$ for expectations and marginals under it, and
$\mathcal{L}=f\cdot\nabla + a\Delta$ for its generator. Given a terminal reward
$r:\mathbb{R}^{d}\to\mathbb{R}$, the object we wish to sample from is the
\emph{reward-tilted} path measure
\begin{equation}
  \frac{\mathrm{d}\mathbb{P}^{\star}}{\mathrm{d}\mathbb{P}}
  \;=\; \frac{e^{r(X_1)}}{Z},
  \qquad
  Z \;=\; \E_{\Base}\!\big[e^{r(X_1)}\big],
  \label{eq:tilt}
\end{equation}
whose terminal marginal is $p^{\star}(x)\propto p_{\Base}(x)\,e^{r(x)}$. This is
the common form of reward-guided generation and of diffusion fine-tuning
\citep{song2021sde,black2024ddpo,uehara2024finetuning}, and it also covers Boltzmann sampling: to target
$\pi\propto e^{-E}$ one sets $r = -E - \log p_{\Base}$, which is how we
instantiate the GMM-40 benchmark in \S\ref{sec:07_experiments}.

Nothing below depends on $f$ having a particular form: if $\mathbb{P}$ is a
pretrained diffusion or flow model, $f$ is its learned drift and the whole
development applies unchanged.\footnote{Note the time convention: $t$ runs from
a deterministic start $X_0=0$ to the model's samples at $t=1$, as in the
diffusion-sampler literature \citep{zhang2021pis,vargas2023dds}, rather than
the denoising direction of DDPM-style presentations.}

\paragraph{The base family we use.}
\label{subsec:base-family}
Our base processes are \emph{pinned interpolants}
\citep{albergo2023interpolants,lipman2023flowmatching}: $X_0=0$, a prescribed
terminal marginal $\nu$, and a Brownian bridge between the two given $X_1=x_1$.
The induced drift $f(x,t)=\E[(X_1-X_t)/(1-t)\mid X_t=x]$ vanishes exactly when
$\nu=\mathcal{N}(0,2aI)$; taking $\nu$ to be a Gaussian mixture instead gives a
base model with nontrivial but analytically known $f$, $\Vf$ and $\log Z$,
which is what makes the study of \S\ref{sec:07_experiments} possible. Details
in App.~\ref{sec:C_experimental_details}.

% ---------------------------------------------------------------------------
\subsection{Control formulation and the value function}

Realizing \eqref{eq:tilt} means adding a drift. For a control $u$ giving the
controlled dynamics $\mathrm{d}X_t=(f+u)\,\mathrm{d}t+\sqrt{2a}\,\mathrm{d}W_t$
with path measure $\mathbb{P}^{u}$, Girsanov's theorem gives
$\mathrm{KL}(\mathbb{P}^{u}\,\|\,\mathbb{P})=\E_{\mathbb{P}^u}\!\int_0^1
\tfrac{1}{4a}\|u_t\|^{2}\,\mathrm{d}t$, so \eqref{eq:tilt} is the solution of
the entropy-regularized control problem \citep{daipra1991stochastic,leonard2014survey}
\begin{equation}
  \max_{u}\ \ \E_{\mathbb{P}^{u}}\!\left[\,r(X_1)
    - \int_0^1 \frac{\|u_t\|^{2}}{4a}\,\mathrm{d}t\right],
  \label{eq:control}
\end{equation}
whose optimum is $\log Z$. The central object is the \emph{soft value function}, the continuous-time
counterpart of the soft value of maximum-entropy reinforcement learning
\citep{haarnoja2017softq,levine2018rlasinference},
\begin{equation}
  \Vf(x,t) \;=\; \log \E_{\Base}\!\big[\,e^{r(X_1)} \,\big|\, X_t=x\,\big],
  \label{eq:value}
\end{equation}
a conditional log-partition function. It solves the HJB equation
$\partial_t \Vf + f\cdot\nabla \Vf + a\Delta \Vf + a\|\nabla \Vf\|^{2}=0$ with
terminal condition $\Vf(\cdot,1)=r$, the optimal control is
\begin{equation}
  u^{\star}(x,t) \;=\; 2a\,\nabla_x \Vf(x,t),
  \label{eq:optimal-control}
\end{equation}
and, since $X_0=0$ is deterministic, $\Vf(0,0)=\log Z$ --- the normalizing
constant, which is the headline metric on external sampler benchmarks
(\S\ref{sec:07_experiments}). Learning $\Vf$ therefore solves the sampling
problem, the fine-tuning problem, and the free-energy-estimation problem at
once, and everything downstream reduces to \emph{regressing $\Vf$ from
Monte-Carlo targets}.

% ---------------------------------------------------------------------------
\subsection{The H-martingale property}
\label{subsec:hmartingale}

Exponentiating \eqref{eq:value} linearizes the HJB equation --- the Hopf--Cole
change of variables underlying linearly-solvable and path-integral control
\citep{kappen2005path,todorov2009efficient}. Writing
$h = e^{\Vf}$, Feynman--Kac gives $\partial_t h + \mathcal{L}h = 0$: $h$ is
space-time harmonic for the base process, so $h(X_t,t)$ is a $\mathbb{P}$-martingale,
and \eqref{eq:optimal-control} is exactly the Doob $h$-transform drift
$f + 2a\nabla\log h$. Concretely, for any $0\le t<s\le 1$,
\begin{equation}
  \boxed{\ e^{\Vf(x,t)} \;=\; \E_{\Base}\!\big[\,e^{\Vf(X_s,\,s)}\,\big|\,X_t=x\,\big]\ }
  \label{eq:hmart}
\end{equation}
with the terminal case $\Vf(\cdot,1)=r$. We call \eqref{eq:hmart} the
H-martingale property; it is the source of every target in this paper, and two
of its features drive the design.

\paragraph{(i) Targets are unbiased in exponential space --- and only there.}
Equation \eqref{eq:hmart} says that a \emph{single} draw of $X_s$ from the base
kernel supplies a scalar $q = \Vf(X_s,s)$ (or $q=r(X_1)$ when $s=1$) satisfying
\begin{equation}
  \E\big[e^{q}\,\big|\,X_t=x\big] \;=\; e^{\Vf(x,t)}.
  \label{eq:unbiased-target}
\end{equation}
The estimator is unbiased for $e^{\Vf}$, not for $\Vf$: the natural log-space
plug-ins are both biased low by a Jensen gap (\S\ref{sec:03_expmse}). This is a
property of the log-partition function rather than of a particular scheme, and
it is what forces the regression into exponential space --- and hence into the
conditioning problem of \S\ref{sec:03_expmse}--\S\ref{sec:04_bregman}.

\paragraph{(ii) The identity is twist-independent.}
Suppose we do not want to draw $X_s$ from the base kernel --- for instance
because base samples rarely reach high-reward regions: importance sampling for
$\E[e^{r}]$ needs on the order of $e^{\mathrm{KL}(\mathbb{P}^{\star}\|\mathbb{P})}$
samples \citep{chatterjee2018sample}, which is the quantitative form of the
sparsity argument of \S\ref{sec:06_sparsity_onpolicy}. Propose instead under any adapted twist
$u$, giving the drift $f+u$ and path measure $\mathbb{P}^{u}$, and correct by
the Radon--Nikodym derivative of the base law with respect to it. By Girsanov's
theorem, restricted to the interval $[t,s]$,
\begin{equation}
  \rho_{t\to s} \;=\; \frac{\mathrm{d}\mathbb{P}}{\mathrm{d}\mathbb{P}^{u}}
  \bigg|_{[t,s]}
  \;=\; \exp\!\left(
    -\frac{1}{2a}\int_t^s u_\tau\cdot \mathrm{d}B_\tau
    \;-\; \frac{1}{4a}\int_t^s \|u_\tau\|^{2}\,\mathrm{d}\tau
  \right),
  \label{eq:girsanov-rho}
\end{equation}
where $\mathrm{d}B_\tau=\sqrt{2a}\,\mathrm{d}W^{u}_\tau$ is the driving noise
under $\mathbb{P}^{u}$ (so the stochastic integral is an It\^o integral against
the noise actually realized by the sampler). Under Novikov's condition
$\rho_{t\to s}$ is an exponential martingale with unit mean, and \eqref{eq:hmart}
becomes
\begin{equation}
  e^{\Vf(x,t)} \;=\; \E_{\mathbb{P}^{u}}\!\big[\,\rho_{t\to s}\, e^{\Vf(X_s,s)}\,\big|\,X_t=x\,\big].
  \label{eq:girsanov}
\end{equation}
In an Euler--Maruyama discretization \eqref{eq:girsanov-rho} factorizes over
steps, each contributing
$-\,u\cdot\Delta B/(2a) - \|u\|^{2}\Delta t/(4a)$ with
$\Delta B\sim\mathcal{N}(0,2a\Delta t\,I)$; this per-step ratio is the
incremental weight of twisted-SMC samplers, a construction that predates its
diffusion applications
\citep{whiteley2014twisted,guarniero2017iapf,heng2020csmc,singhal2025fksteering}. The twist $u$ is therefore \emph{free}: it
changes the variance of the target, never its mean. The proposal is a
variance-reduction device, not a modeling choice --- which is what licenses the
importance-sampling and SMC machinery of \S\ref{sec:06_sparsity_onpolicy}
without changing the object being learned.

% ---------------------------------------------------------------------------
\subsection{How training pairs are formed}
\label{subsec:targets}

A learner $\Vf_\theta$ is fit to pairs $(x_t, q)$, of two kinds. \emph{Anchored} targets terminate at the true reward and satisfy
\eqref{eq:unbiased-target} exactly, so $e^{q}$ is unbiased for $e^{\Vf(x_t,t)}$
however wrong the current $\Vf_\theta$ is. \emph{Bootstrapped} targets
substitute the learner's own value at a later time, and satisfy no such
identity. Both appear below; the distinction matters, and we are explicit about
which guarantee is in force.

\paragraph{Off-policy (bridge) anchors.}
Draw a terminal point $x_1\sim\nu$ from the base terminal marginal, place $x_t$
on the base bridge to it, and take $q=r(x_1)$. For the pinned interpolants of
\S\ref{sec:02_setting} that bridge is the Brownian bridge in closed form,
\begin{equation}
  x_t \;=\; t\,x_1 + \sqrt{2a\,t(1-t)}\;\varepsilon,
  \qquad \varepsilon\sim\mathcal{N}(0,I),
  \label{eq:fwd-bridge}
\end{equation}
so the sampled joint $(x_t,x_1)$ is \emph{exactly} the base joint and $q$ is
anchored in the sense above. What matters is that $x_1$ is drawn from the base
marginal $\nu$: anchoring on a reward-aware or held-out data distribution
instead breaks the conditional and biases the targets.

These anchors are cheap \emph{per sample} --- no trajectory is integrated, and
each pair costs a single evaluation of $r$ --- and they cover the whole
marginal. Two caveats. First, $x_1$ is drawn without regard to $r$, which is
what fails in high dimension (\S\ref{sec:06_sparsity_onpolicy}). Second, cheap
per sample is not the same as cheap in total: because the anchors are
uninformative about where $r$ is large, reaching a given accuracy can take many
more of them, and hence many more reward evaluations, than a targeted scheme.
When $r$ is itself expensive --- a quantum-chemistry calculation, a docking
score, a simulator rollout --- total calls to $r$, not wall-clock per sample,
is the budget that binds, and the ranking between the constructions can invert.
Our benchmarks all have cheap analytic rewards, so we report sample counts
rather than reward-call counts; we flag this as a limitation in
\S\ref{sec:09_conclusion}.

\paragraph{On-policy (bootstrapped) targets, and backward-noising.}
The other two constructions are used by the experiments but are not needed to
follow the argument, so we state them briefly and defer the details to
App.~\ref{sec:B_onpolicy_zoo}. A twisted SMC sweep supplies \emph{bootstrapped}
targets, formed by aggregating the learner's own values at the next step; these
satisfy no exact identity, being unbiased in exponential space for the one-step
backup of the current $\Vf_\theta$ rather than for $\Vf$, and their usefulness
rests on the terminal step carrying $q=r(X_1)$ exactly. Separately, because
$e^{\Vf}$ is harmonic, a target attached at $(x_s,s)$ can be reattached at a
point noised backwards from it through the base backward kernel
\begin{equation}
  X_t \,\big|\, X_s = x_s \;\sim\; \mathcal{N}\!\Big(\tfrac{t}{s}\,x_s,\;
    2a\,\tfrac{t(s-t)}{s}\,I\Big), \qquad t<s,
  \label{eq:backward-kernel}
\end{equation}
which for the pinned interpolants does not depend on the terminal law $\nu$ at
all (App.~\ref{sec:D_proofs}), so the construction is available for base models
with nontrivial drift. It is exact when the noised point is base-distributed
and approximate under a twist (\S\ref{sec:06_sparsity_onpolicy}).

\medskip
\noindent In every case the learning problem has the same shape: given pairs
$(x_t,q)$ with $\E[e^{q}\mid x_t] = e^{\Vf(x_t,t)}$, fit $p=\Vf_\theta(x_t)$.
The rest of the paper concerns the two remaining degrees of freedom --- the
divergence applied to $(p,q)$, and the sampler that produces the pairs.
